\documentclass[10pt,a4paper]{amsart}
\usepackage[margin=19mm]{geometry}
\usepackage{amsmath,amssymb,mathtools}
\usepackage[scr=rsfs,cal=euler]{mathalfa}
\usepackage{xfrac}
\usepackage[unicode,hidelinks,bookmarks=false]{hyperref}
\newcommand{\ie}{i.e.\ }
\newcommand{\cf}{cf.\ }
\newcommand{\resp}{respectively\ }
\newcommand{\Subsection}{\subsection}
\providecommand{\nobreakdash}{\nobreak\hbox{-}}
\setlength{\emergencystretch}{2em}
\multlinegap0pt
\usepackage{bm}
\usepackage{amssymb}
\newtheorem{lemma}{Lemma}
\newtheorem{theorem}{Theorem}
\title{On the $\ell_p$ and $\ell_{p,q}$ norms of Cauchy--Toeplitz matrices}
\author{Tserendorj Batbold}
\date{Source: arXiv:2501.17566v1 (CC BY 4.0)}
\begin{document}
\maketitle

\noindent\emph{Source and licence.} On the $\ell_p$ and $\ell_{p,q}$ norms of Cauchy--Toeplitz matrices. Version arXiv:2501.17566v1 (CC BY 4.0), distributed by arXiv under the Creative Commons Attribution 4.0 International licence, http://creativecommons.org/licenses/by/4.0/, as recorded in the arXiv licence metadata for this exact version and shown on the abstract page https://arxiv.org/abs/2501.17566v1. Full licence notice: \texttt{licenses/arxiv-2501.17566-CC-BY-4.0.txt}.\par\smallskip

\noindent\textit{Editorial scope.} Counted unit: the article's theorem mainthm1 (source0.tex lines 307--317). The article's complete original proof of it is delivered with it (source0.tex lines 318--408, one \texttt{proof} environment of 91 lines). No mathematics is written, repaired or re-derived: the statement and the proof are the article's own lines, in the article's own notation, and no retained line is substituted or re-flowed.  Retained uncounted as the source-local context the proof needs: lines 34--36; lines 89--93; lines 210--212.  Nothing was omitted from any retained span, so the package carries no omission record.  No retained line contains a citation, so the wrapper carries no bibliography and no bibliography entry.  No retained line contains a figure environment, an image-inclusion command or drawing code, so no image is delivered and the package declares no asset dependency.  Editorial formatting only, in addition: an AMS \texttt{amsart} preamble that restates the article's own declarations and macros used by the retained lines, namely the environments \texttt{lemma}, \texttt{theorem} on separate counters, together with the standard packages those lines need.  The retained environments print in this document as Lemma 1, Theorem 1 in order of appearance, whereas the article prints the same items with the numbering of its own full document.  The complete native source of the article as distributed in the arXiv e-print for this exact version is bundled unchanged as \texttt{sources/172371/source0.tex}.
\begin{equation}\label{matrix}
   T_n=\left[\frac{2}{1+2(i-j)}\right]_{i, j=1}^n.
\end{equation}

\begin{lemma}\label{lem2}
    Let $p\geq1$. Then the following inequality
    $$3-2^p+\frac{1}{3^p}-\left(\frac{2}{3}\right)^p>0$$
    holds for $1\leq p<\delta$. The opposite inequality holds $p>\delta$. Here $\delta=1.40485...$ is the unique root of $2^\delta + 6^\delta = 3^{\delta + 1} + 1$.
\end{lemma}

\begin{equation}\label{generalineq1}
   2\leq n^{-\frac{1}{p}}\|T_n\|_{p}\leq m^{-\frac{1}{p}}\|T_m\|_{p}.
\end{equation}

\begin{theorem}\label{mainthm1}
Let $T_n$ be as in (\ref{matrix}), and let $1\leq p<q\leq\infty$. Then the following inequality
\begin{equation}\label{mainthmeq1}
n^{-\frac{1}{q}}\left\|T_{n}\right\|_{p,q}\geq 4\left(\frac{1}{2^{p}-1}\right)^{\frac{1}{p}}, 
\end{equation}
holds for $n\geq3$ and for $n=2, (p\geq\delta~~\text{or}~~1<p<\delta, q\geq \delta_p)$. The opposite inequality holds for $n=1$ and for $n=2, (p=1~~\text{or}~~1< p<\delta, q<\delta_p )$. Here $q=\delta_p$ is the unique root of the following equation
\begin{equation}\label{delta}
2^{-p}\left(\left(1+\frac{1}{3^p}\right)^\frac{q}{p}+2^\frac{q}{p}\right)^p=2^{pq}\left(\frac{1}{2^p-1}\right)^q,
\end{equation}
for fixed $p$. 
\end{theorem}

\begin{proof} {\bf Case $\boldsymbol{(n\geq3)}$}: By the power mean inequality, we have
\begin{align}\label{monoeq1}
n^{-\frac{1}{q}}\left\|T_{n}\right\|_{p,q}&=n^{-\frac{1}{q}}\left[\sum_{j=1}^n\left(\sum_{i=1}^n|a_{ij}|^p\right)^\frac{q}{p}\right]^\frac{1}{q}=\left[\frac{1}{n}\sum_{j=1}^n\left(\sum_{i=1}^n|a_{ij}|^p\right)^\frac{q}{p}\right]^\frac{1}{q}\notag\\
&=M_q\left(\underline{\left(\sum_{i=1}^n|a_{ij}|^p\right)^\frac{1}{p}}\right)\notag\\
&>M_p\left(\underline{\left(\sum_{i=1}^n|a_{ij}|^p\right)^\frac{1}{p}}\right)=\left[\frac{1}{n}\sum_{j=1}^n\left(\sum_{i=1}^n|a_{ij}|^p\right)^\frac{p}{p}\right]^\frac{1}{p}\notag\\
&=n^{-\frac{1}{p}}\left[\sum_{j=1}^n\sum_{i=1}^n|a_{ij}|^p\right]^\frac{1}{p}=n^{-\frac{1}{p}}\|T_n\|_p.
\end{align}  
Applying  (\ref{generalineq1}), gives
\begin{equation}\label{mainthm1eq2}
n^{-\frac{1}{p}}\|T_n\|_p>3^{-\frac{1}{p}}\left\|T_{3}\right\|_{p}=3^{-\frac{1}{p}}\left(2^{p}\left(5+\frac{1}{3^{p-1}}+\frac{1}{5^{p}}\right)\right)^{\frac{1}{p}}.
\end{equation}
Combining (\ref{monoeq1}) and (\ref{mainthm1eq2}) leads to
$$
n^{-\frac{1}{q}}\left\|T_{n}\right\|_{p,q} >3^{-\frac{1}{p}}\left(2^{p}\left(5+\frac{1}{3^{p-1}}+\frac{1}{5^{p}}\right)\right)^{\frac{1}{p}}.
$$
To prove the inequality (\ref{mainthmeq1}) for $n\geq3$, it suffices to show that the following inequality
$$
3^{-\frac{1}{p}}\left(2^{p}\left(5+\frac{1}{3^{p-1}}+\frac{1}{5^{p}}\right)\right)^{\frac{1}{p}}>4\left(\frac{1}{2^{p}-1}\right)^{\frac{1}{p}},
$$
i.e.,
$$
2^{p+1}-5+3\left(\frac{2}{3}\right)^{p}+\left(\frac{2}{5}\right)^{p}-\frac{1}{3^{p-1}}-\frac{1}{5^p}>0.
$$
Considering the function $h$ on $[1,+\infty)$ defined by 
$$
h(p)=2^{p+1}-5+3\left(\frac{2}{3}\right)^{p}+\left(\frac{2}{5}\right)^{p}-\frac{1}{3^{p-1}}-\frac{1}{5^p}.
$$
Differentiation yields
\begin{align*}
h'(p)&=2^{p+1} \ln 2+3\left(\frac{2}{3}\right)^{p} \ln \frac{2}{3}+\left(\frac{2}{5}\right)^{p} \ln \frac{2}{5}+\frac{1}{3^{p-1}}\ln3+\frac{1}{5^p}\ln 5\\
&>2^{p+1} \ln 2+3\left(\frac{2}{3}\right)^{p} \ln \frac{2}{3}+\left(\frac{2}{5}\right)^{p} \ln \frac{2}{5}:=h_1(p)
\end{align*}
and 
$$
h_1'(p)=  2^{p+1} (\ln 2)^2+3\left(\frac{2}{3}\right)^{p} \left(\ln \frac{2}{3}\right)^2+\left(\frac{2}{5}\right)^{p} \left(\ln \frac{2}{5}\right)^2>0.
$$
Thus, we deduce that for $p\geq1$,
 $$
h_1(p) \geq h_1(1)=4\ln 2+2\ln \frac{2}{3}+\frac{2}{5}\ln \frac{2}{5}\approx1.5951\ldots>0.
$$
Therefore, the function $h(p)$ is increasing on $[1,+\infty[$,  which implies that 
$$h(p)>h(1)=\frac{1}{5}>0$$
for $p\geq1$.

{\bf Case $\boldsymbol{(n=2, p\geq\delta)}$:}  From the inequality (\ref{monoeq1}), we obtain
$$
2^{-\frac{1}{q}}\left\|T_{2}\right\|_{p,q}>2^{-\frac{1}{p}}\left\|T_{2}\right\|_{p}.
$$
Clearly, it suffices to show that 
\begin{equation}\label{eq12}
2^{-\frac{1}{p}}\left\|T_{2}\right\|_{p} \geq 4\left(\frac{1}{2^{p}-1}\right)^{\frac{1}{p}},
\end{equation}
This inequality is equivalent to
 $$3-2^p+\frac{1}{3^p}-\left(\frac{2}{3}\right)^p\leq0,$$
which is true by the Lemma \ref{lem2}.

 {\bf Case $\boldsymbol{(n=2, 1\leq p<\delta)}$:} If $p=1$, then by the power mean inequality, we have
$$2^{-\frac{1}{q}}\left\|T_{2}\right\|_{1, q}=2^{1-\frac{1}{q}}\left(\left(\frac{4}{3}\right)^q+2^q\right)^\frac{1}{q}<\lim_{q\to\infty}2^{1-\frac{1}{q}}\left(\left(\frac{4}{3}\right)^q+2^q\right)^\frac{1}{q}=4.$$
In other words, the inequality (\ref{mainthmeq1}) does not hold for $p=1$.
 
Similarly to the previous case, one has for $1\leq p<\delta$,
\begin{equation}\label{thmeqq2}
2^{-\frac{1}{p}}\left\|T_{2}\right\|_{p} < 4\left(\frac{1}{2^{p}-1}\right)^{\frac{1}{p}}.
\end{equation}


If $1< p<\delta$, then by the power mean inequality, we have  
\begin{align*}
2^{-\frac{1}{q}}\left\|T_{2}\right\|_{p,q}&=2^{1-\frac{1}{q}}\left(\left(1+\frac{1}{3^p}\right)^\frac{q}{p}+2^\frac{q}{p}\right)^\frac{1}{q}\\
&<\lim_{q\to+\infty}2^{1-\frac{1}{q}}\left(\left(1+\frac{1}{3^p}\right)^\frac{q}{p}+2^\frac{q}{p}\right)^\frac{1}{q}=2^{1+\frac{1}{p}}.
\end{align*}
It is easy to see that
$$4\left(\frac{1}{2^{p}-1}\right)^{\frac{1}{p}}<2^{1+\frac{1}{p}}.$$
By the continuity and monotonicity of power mean, there are $q'$ and $q''$ such that $p<q'\leq q''$ and 
$$
2^{-\frac{1}{p}}\left\|T_{2}\right\|_{p}<2^{-\frac{1}{q}}\left\|T_{2}\right\|_{p, q'}\leq4\left(\frac{1}{2^p-1}\right)^{\frac{1}{p}}\leq2^{-\frac{1}{q}}\left\|T_{2}\right\|_{p, q''}<2^{1+\frac{1}{p}}.
$$
Hence, there is exist a unique $q=\delta_p$ for fixed $p$, such that
$$2^{-\frac{1}{q}}\left\|T_{2}\right\|_{p,q}=4\left(\frac{1}{2^p-1}\right)^{\frac{1}{p}},$$
that is,
$$2^{-p}\left(\left(1+\frac{1}{3^p}\right)^\frac{\delta_p}{p}+2^\frac{\delta_p}{p}\right)^p=2^{p\delta_p}\left(\frac{1}{2^p-1}\right)^{\delta_p}.$$
Therefore, by the power mean inequality, we have
$$
2^{-\frac{1}{q}}\left\|T_{2}\right\|_{p,q}<4\left(\frac{1}{2^{p}-1}\right)^{\frac{1}{p}},
$$
for $1<q<\delta_p$, and 
$$
2^{-\frac{1}{q}}\left\|T_{2}\right\|_{p,q}\geq4\left(\frac{1}{2^{p}-1}\right)^{\frac{1}{p}},
$$
for $q\geq\delta_p$. We thus complete the proof.
\end{proof}

\end{document}
